% =====================================================================
%  3.2 Andregradsfunksjoner — fasit
% =====================================================================

\fk{Del 1 — Graf, nullpunkter og symmetrilinje}
\fo{1.1} a) opp, $(0, 4)$ \ b) ned, $(0, 7)$ \ c) ned, $(0, 3)$
\fo{1.2} a) $x = 4$ \ b) $x = -\frac{3}{2}$ \ c) $x = 5$
\fo{1.3} a) $x = -1$ og $x = 5$ \ b) $x = \frac{-1 + 5}{2} = 2$ \ c) $(2, f(2)) = (2, -9)$
\fo{1.4} a) Nullpunkt $x = -5$ og $x = 3$. Symmetrilinje $x = -1$. Bunnpunkt $(-1, -16)$.
b) Nullpunkt $x = 2$ og $x = 4$. Symmetrilinje $x = 3$. Toppunkt $(3, 1)$.
c) $D = -24 < 0$, ingen nullpunkter. Symmetrilinje $x = 1$. Bunnpunkt $(1, 3)$.
\fo{1.5} Feilen er overgangen til $-12 - 24 + 5$: $3 \cdot (-2)^2 = 3 \cdot 4 = 12$, ikke $-12$.
Riktig: $f(-2) = 12 - 24 + 5 = -7$. Siden $a = 3 > 0$ er det et bunnpunkt, ikke et toppunkt: $(-2, -7)$.
\fo{1.6} a) Alle har samme form og samme symmetrilinje. De er flyttet opp eller ned i forhold til hverandre.
b) Symmetrilinje $x = -3$. Bunnpunkt $(-3, c - 9)$. Bunnpunktene er $(-3, -11)$, $(-3, -6)$ og $(-3, -1)$.
\fo{1.7} a) Noen ganger. $x^2 + 1$ har ingen nullpunkter, men $x^2 - 1$ har to.
b) Noen ganger. For $x^2 + 4$ er bunnpunktet $(0, 4)$ og $y = c$. For $x^2 + 4x + 4$ er bunnpunktet $(-2, 0)$, og $c = 4$. Det stemmer bare når $b = 0$.
c) Aldri. Når $a > 0$ vender parabelen opp og har et bunnpunkt.
\fo{1.8} a) $-\frac{b}{2} = -1$ gir $b = 2$. Nullpunktet gir $4 + 4 + c = 0$, så $c = -8$. Altså $f(x) = x^2 + 2x - 8$.
b) Symmetrilinjen ligger midt mellom nullpunktene. Fra $x = 2$ til $x = -1$ er avstanden 3. Det andre nullpunktet ligger like langt på den andre siden: $-1 - 3 = -4$.
\fo{Sjekk 1} a) $x = 5$ \ b) ned, $(0, 4)$ \ c) Summen av nullpunktene er $-\frac{b}{a}$. Gjennomsnittet er $-\frac{b}{2a}$, og det er midtpunktet mellom dem.

\fk{Del 2 — Fullstendige kvadraters metode og toppunktform}
\fo{2.1} a) $(x + 1)^2 + 4$ \ b) $(x - 5)^2 + 5$ \ c) $(x - 4)^2 - 7$
\fo{2.2} a) Bunnpunkt $(2, 7)$, $V_f = [7, \rightarrow\rangle$ \ b) Toppunkt $(-1, 4)$, $V_g = \langle\leftarrow, 4]$ \ c) Bunnpunkt $(-4, -2)$, $V_h = [-2, \rightarrow\rangle$
\fo{2.3} $f(x) = x^2 + 16x + 64 - 64 + 60 = (x + 8)^2 - 4$. Bunnpunkt $(-8, -4)$.
\fo{2.4} a) $2(x + 3)^2 - 11$, bunnpunkt $(-3, -11)$ \ b) $3(x - 1)^2 + 2$, bunnpunkt $(1, 2)$ \ c) $-(x - 2)^2 + 5$, toppunkt $(2, 5)$
\fo{2.5} Feilen er siste overgang: $-18 + 5 = -13$, ikke $-23$. Riktig: $f(x) = 2(x - 3)^2 - 13$, bunnpunkt $(3, -13)$.
\fo{2.6} a) $h(t) = -5(t - 1)^2 + 7$ \ b) Toppunktet er $(1, 7)$: ballen når 7 m etter 1 s. \ c) $h(0) = 2$, altså 2 m.
\fo{2.7} a) Toppunktform $f(x) = a(x - 2)^2 - 3$. Punktet $(0, 5)$ gir $4a - 3 = 5$, så $a = 2$. Altså $f(x) = 2(x - 2)^2 - 3 = 2x^2 - 8x + 5$.
b) To nullpunkter. Parabelen vender opp ($a = 2 > 0$) og bunnpunktet $(2, -3)$ ligger under $x$-aksen. Da må grafen krysse aksen på begge sider av $x = 2$.
\fo{2.8} a) Lengden langs veggen er $40 - 2x$, fordi to sider er $x$. Da er $A(x) = x(40 - 2x) = -2x^2 + 40x$.
b) $A(x) = -2(x - 10)^2 + 200$. Største areal er $200\text{ m}^2$ for $x = 10$. Hagen er 10 m $\times$ 20 m.
c) $A(25) = -250$ er negativt. Lengden $40 - 2 \cdot 25 = -10$ kan ikke være negativ. Bredden må være mellom 0 og 20.
\fo{2.9} a) Aldri. Toppunktet $(p, q)$ ligger under $x$-aksen, og parabelen vender ned. Da er $f(x) \le q < 0$ for alle $x$.
b) Noen ganger. $2(x - 1)^2 + 3$ har ingen nullpunkter. $-2(x - 1)^2 + 3$ har to.
c) Noen ganger. $x^2 - 4x + 1$ og $2x^2 - 8x + 5$ har begge bunnpunkt $(2, -3)$, men er ulike funksjoner. Det kreves også samme $a$.
\fo{Sjekk 2} a) $(x - 6)^2 + 4$ \ b) Bunnpunkt $(-1, -6)$, $V_g = [-6, \rightarrow\rangle$ \ c) Fordi $(x - p)^2 \ge 0$ og $a > 0$, er $a(x - p)^2 \ge 0$. Da er $f(x) = a(x-p)^2 + q \ge q$.

\fk{Ekstra utfordring}
\fo{U1} $f(x) = a(x + 1)(x - 5)$. Punktet $(0, -10)$ gir $-5a = -10$, så $a = 2$. Altså $f(x) = 2x^2 - 8x - 10$. Symmetrilinje $x = \frac{-1 + 5}{2} = 2$, $f(2) = -18$. Toppunktformen er $2(x - 2)^2 - 18$, så bunnpunktet er $(2, -18)$ (parabelen vender opp, så det er et bunnpunkt).
\fo{U2} a) $I(p) = p(400 - 20p) = -20p^2 + 400p = -20(p - 10)^2 + 2000$. Pris 10 kr gir størst inntekt, 2000 kr.
b) Salget $400 - 20p$ blir negativt når $p > 20$. Modellen gjelder bare for $0 \le p \le 20$. Negativ inntekt gir ingen mening.
\fo{U3} a) Programmet skriver ut \texttt{-3.0}. Toppunktformen er $(x - 2)^2 - 3$, så minimum er $-3$ for $x = 2$. Verdien $x = 2$ er et steg i programmet.
b) Programmet skriver ut \texttt{-0.005}. Fra toppunktformen er minimum $-\frac{9}{1600} = -0{,}005625$ for $x = 0{,}075$. Det punktet ligger ikke på et steg. Nærmeste steg er $x = 0{,}1$, som gir $-0{,}005$.
\fo{U4} $f(x) = (x - 1)^2 + c - 1$ med bunnpunkt $(1, c - 1)$. Parabelen vender opp.
$c < 1$: bunnpunktet ligger under aksen, og det er to nullpunkter. $c = 1$: bunnpunktet ligger på aksen, og det er ett nullpunkt. $c > 1$: bunnpunktet ligger over aksen, og det er ingen nullpunkter.
\fo{U5} $a\left(x + \frac{b}{2a}\right)^2 = a\left(x^2 + \frac{b}{a}x + \frac{b^2}{4a^2}\right) = ax^2 + bx + \frac{b^2}{4a}$. Trekk fra $\frac{b^2}{4a}$ og legg til $c$: det gir $ax^2 + bx + c$.
a) Symmetrilinjen er $x = -\frac{b}{2a}$ og $q = c - \frac{b^2}{4a}$.
b) $q = \frac{4ac - b^2}{4a} = -\frac{D}{4a}$. Er $D < 0$ og $a > 0$, er $q > 0$. Bunnpunktet ligger over $x$-aksen, så det er ingen nullpunkter.
\fo{U6} a) $b = -4$: $(2, -4)$. $b = -2$: $(1, -1)$. $b = 0$: $(0, 0)$. $b = 2$: $(-1, -1)$. $b = 4$: $(-2, -4)$.
b) Punktene ligger på $y = -x^2$. Bevis: $x^2 + bx = (x + \frac{b}{2})^2 - \frac{b^2}{4}$, så bunnpunktet er $(-\frac{b}{2}, -\frac{b^2}{4})$. Med $x = -\frac{b}{2}$ er $-x^2 = -\frac{b^2}{4} = y$.
